\documentclass[11pt]{article}

\usepackage[margin=1in]{geometry}
\usepackage{amsmath,amssymb,amsfonts,amsthm}
\usepackage{booktabs}
\usepackage{multirow}
\usepackage{graphicx}
\usepackage{hyperref}
\usepackage{xcolor}
\hypersetup{
    colorlinks=true,
    linkcolor=black,
    citecolor=black,
    urlcolor=blue!50!black,
}
\providecommand{\doi}[1]{\href{https://doi.org/#1}{doi:#1}}
\usepackage[expansion=false]{microtype}
\usepackage{authblk}
\usepackage{natbib}

\newcommand{\GROUND}{\Gamma_{\mathrm{i}}}
\newcommand{\KL}{D_{\mathrm{KL}}}
\newcommand{\G}{g}

\newtheorem{theorem}{Theorem}
\newtheorem{proposition}{Proposition}
\newtheorem{lemma}{Lemma}
\theoremstyle{definition}
\newtheorem{definition}{Definition}
\theoremstyle{remark}
\newtheorem*{remark}{Remark}

\title{The Intrinsic GROUND Ranker with Volatility-Divergence Discount: \\
       SP100 Weekly Credit Spreads (2020--2026)}

\author{Peter J.~Mercurio, PhD}
\affil{\small \texttt{peter.mercurio@alumni.utoronto.ca}}
\date{June 2026}


\begin{document}
\maketitle

\begin{abstract}
We present a six-and-a-half-year empirical validation of the GROUND ranking
framework applied to weekly vertical credit spreads on the SP100 universe.
We adopt the intrinsic GROUND ratio in its \emph{Kelly-EV-with-entropic-discount}
multiplicative form
\[
\GROUND(p, x, k) \;=\; E(p, x) \cdot \exp\!\bigl(-k \cdot \KL(P \,\|\, Q)\bigr),
\]
where $E(p, x) := M(p, x) - 1$ is the Kelly EV computed at the closed-form
Kelly-optimal sizing $w^\star$, $M = \exp(\ell(w^\star))$ is the per-trade
wealth multiplier under classical log-utility, and the entropic discount
$\exp(-k \cdot \KL) \in (0, 1]$ penalizes per-spread divergence between two
risk-neutral probability distributions $P$ and $Q$ constructed from
\emph{realized} and \emph{implied} volatility respectively via the
Black--Scholes $d_2$. We show that this choice of reference --- realized
versus implied volatility through the BS $d_2$ --- captures a per-spread
variance-risk-premium signal that empirical-bucket-conditioned references
do not. The amplification parameter $k = 10$ and the qualification
threshold $0.075$ are selected from a two-dimensional sweep on the full
sample. On the 2020--2025 window
(\textbf{1{,}111 trades}, 297 weeks, qty=2, post-hoc partial-WIN haircut
to 50\%~of intrinsic for live pin-risk realism) the strategy posts
\textbf{Sharpe 2.05 weekly} (1.99 daily), \textbf{maximum drawdown 8.0\%},
\textbf{CAGR 33.4\%}, and a bankroll-independent yield-on-wager of
\textbf{28.5\%}. On the strict 2026 holdout
(\textbf{[OOT-PICKS]} trades, [OOT-WEEKS] weeks, never used in any
parameter selection) the strategy posts \textbf{Sharpe [OOT-SHARPE]
weekly}, \textbf{maximum drawdown [OOT-DD]\%}, \textbf{yield [OOT-YIELD]\%}.
SPY buy-and-hold over the same windows posts Sharpe 0.71 and Sharpe
[OOT-SPY-SHARPE]. Single-factor baselines that rank by Kelly~EV alone
posts a per-pick mean of $+\$22.10$ at 44.8\% win-rate; the entropic
discount swaps 32\% of those picks for ones with 51.7\% win-rate,
lifting the per-pick mean to $+\$24.77$ at 52.6\% win-rate without any
change in pick count. Crucially the discount factor itself is empirically
uncorrelated with realized P\&L (Pearson $+0.004$ across 22{,}702
pre-filter candidates): the divergence does not predict outcomes
directly, it acts as a tail filter that suppresses the candidates whose
implied volatility is signaling regime risk realized vol cannot see.
\end{abstract}


\section{Introduction}

The GROUND ratio was introduced in \citet{mercurio2020} as a portfolio
ranking criterion that simultaneously rewards Kelly growth and penalizes
deviation from a reference distribution by an information-theoretic
divergence. In its original form GROUND was a per-week comparative score
against the lowest-divergence candidate in the menu; in this paper we
adopt the intrinsic v2 form
\begin{equation}
\GROUND(p, x, k) \;=\; E(p, x) \cdot \exp\!\bigl(-k \cdot \KL(P \,\|\, Q)\bigr),
\label{eq:ground}
\end{equation}
in which each candidate is scored independently against its own choice
of reference distribution. The wealth multiplier $M(p, x) =
\exp\!\bigl(\ell(w^\star)\bigr)$ is built from the closed-form Kelly log-growth
$\ell(w^\star; p, x) = E_p[\ln(1 + w^\star X)]$ at the Kelly-optimal
sizing fraction $w^\star$ on per-spread payoffs
$\{+b,\, +\alpha b,\, -1\}$, where $b = \mathrm{credit}/\mathrm{max\_loss}$
is the spread's credit-to-risk ratio and $\alpha = \alpha(b)$ is the
principled partial multiplier derived in \citet{mercurio2020}. The Kelly
EV $E(p, x) := M - 1$ is the variance-adjusted expected wealth gain per
trade at Kelly sizing; log-utility bakes in variance via Jensen so
no separate volatility penalty is required at the level of
$E$ itself.

\paragraph{Connection to Hansen--Sargent multiplier preferences.}
Defining the growth signal $g(p, x) := \ln E(p, x)$ (well-defined on the
$E > 0$ growth-positive cone), taking the natural logarithm of
\eqref{eq:ground} yields
\begin{equation}
\ln \GROUND(p, x, k) \;=\; g(p, x) \;-\; k \cdot \KL(P \,\|\, Q) \;=\; J_k(p, x),
\label{eq:hs}
\end{equation}
where $J_k$ is the multiplier-preferences functional of
\citet{hansen2008robustness} and the variational preferences class of
\citet{maccheroni2006ambiguity}. The multiplicative form
\eqref{eq:ground} and the additive form \eqref{eq:hs} produce identical
selection by monotonicity of $\exp$; the multiplicative form has the
advantage of being directly readable as a percent risk-adjusted
expected return per trade.

\paragraph{The contribution of this paper.}
The original \citet{mercurio2020} paper validated GROUND against a
\emph{uniform} reference distribution $u$ on the three-outcome
partition. Subsequent revisions of the canonical replaced $u$ with an
\emph{empirical bucket-conditioned} reference $P_{\mathrm{hist}}$
estimated from rolling historical realizations of similarly-bucketed
spreads. In this paper we propose and validate a third choice: $P$ is
the BS $d_2$-induced probability distribution computed using a
tenor-matched realized-volatility estimate, while $Q$ is the same BS
$d_2$ distribution computed using the option's market-implied
volatility per leg. This choice is theoretically attractive
(it operationalizes the per-spread variance risk premium directly in
probability space rather than in vol space) and empirically dominant
(it beats the bucket-conditioned reference in the head-to-head
out-of-sample test reported in Section~\ref{sec:rvvsiv-choice}). We
also remove the binary directional regime gate that had been part of
prior canonicals; the empirical justification is given in
Section~\ref{sec:no-regime}.

\paragraph{Roadmap.}
Section~\ref{sec:framework} develops the closed-form Kelly-with-partials
math, the wealth-multiplier construction, and the GROUND-Hansen-Sargent
identity in full. Section~\ref{sec:rvvsiv-choice} introduces and
motivates the \emph{rv\_vs\_iv} divergence choice and gives the head-to-head
comparison against the empirical-bucket reference.
Section~\ref{sec:methodology} describes the universe, instrument
selection, entry credit basis, partial-zone arithmetic, and the
earnings and ex-dividend gates that the live implementation applies
to candidate spreads. Section~\ref{sec:results} reports the headline
in-sample backtest (2020--2025). Section~\ref{sec:oot} reports the
strict 2026 out-of-time validation. Section~\ref{sec:proof} establishes
the GROUND-versus-G-alone improvement empirically. Section~\ref{sec:sensitivity}
gives sensitivity analyses for entry credit basis, the partial-win
haircut, and probability-partition collapses. Section~\ref{sec:conclusion}
concludes.


\section{The intrinsic GROUND framework}
\label{sec:framework}

\subsection{Trade outcomes and payoffs}

A vertical credit spread on an underlying with short-leg strike
$K_{\text{short}}$, long-leg strike $K_{\text{long}}$, and entry credit
$c$ has spread width $W = |K_{\text{short}} - K_{\text{long}}|$ and
maximum loss $L = W - c$. The realized P\&L at expiry as a function of
the underlying settlement price $S_T$ is piecewise linear:
\begin{equation}
\Pi(S_T) \;=\;
\begin{cases}
+c & \text{outside (full credit kept)}\\[2pt]
+c - \mathrm{intrinsic}(S_T) & \text{between strikes (partial)}\\[2pt]
-L & \text{past long strike (max loss)}
\end{cases}
\end{equation}
We rescale into per-unit-at-risk units by dividing by $L$; the spread
becomes a three-outcome bet with payoffs $\{+b, +\alpha b, -1\}$ where
$b := c / L$ is the credit-to-risk ratio and $\alpha = \alpha(b)$ is
the partial multiplier derived below.

\subsection{Partial multiplier $\alpha(b)$}

Assume $S_T$ is approximately uniformly distributed across the partial
zone $(K_{\text{long}}, K_{\text{short}})$ conditional on the spread
expiring in the partial zone at all.\footnote{This is the approximation
used in \citet{mercurio2020}; small in practice for the narrow
single-strike spreads we consider.} The expected partial payoff is the
average of $+c$ (at the short strike, breakeven satisfied with no
intrinsic owed) and $-L$ (at the long strike, full intrinsic owed):
\begin{equation}
E_{\mathrm{partial}}[\Pi] \;=\; \tfrac{1}{2}(+c) + \tfrac{1}{2}(-L)
\;=\; \tfrac{1}{2}(c - L)
\;=\; \tfrac{L}{2}\!\left(\tfrac{c-L}{L}\right)
\;=\; \tfrac{L}{2}(b - 1).
\end{equation}
Dividing by $L$ to get the per-unit-at-risk partial payoff:
$\alpha b = (b - 1)/2$, hence
\begin{equation}
\alpha(b) \;=\; \frac{b - 1}{2b}.
\end{equation}
For the boundary case $b = 1$, $\alpha = 0$, i.e., the average partial
zone pays nothing per unit at risk. For credit-rich spreads ($b > 1$)
$\alpha > 0$; for credit-thin spreads ($b < 1$) $\alpha < 0$, the
partial zone is on average punitive. The 5.5-year empirical
distribution of $\alpha$ in our SP100 sample ranges from $-1.17$
(credit-thin, deep-OTM) to $-0.13$ (typical at-the-money).

\subsection{Kelly sizing on a three-outcome bet}

For payoffs $\{+b, +\alpha b, -1\}$ under probability vector
$(p, r_o, q)$ with $p + r_o + q = 1$, the Kelly log-growth at sizing
fraction $w \in (0, 1)$ is
\begin{equation}
\ell(w) \;=\; p \ln(1 + wb) + r_o \ln(1 + w \alpha b) + q \ln(1 - w).
\end{equation}
Differentiating and setting $\ell'(w) = 0$:
\begin{equation}
\frac{p b}{1 + wb} + \frac{r_o \alpha b}{1 + w \alpha b} - \frac{q}{1 - w} \;=\; 0.
\label{eq:kelly-foc}
\end{equation}
Clearing denominators yields a quadratic in $w$ with positive root
$w^\star \in (0, 1)$ for any $(p, r_o, q)$ with positive expected
return. The closed-form Kelly EV is
\begin{equation}
E(p, x) \;=\; \exp\bigl(\ell(w^\star)\bigr) - 1.
\label{eq:E}
\end{equation}

\begin{proposition}[Existence and uniqueness of $w^\star$]
\label{prop:wstar}
For every probability vector $(p, r_o, q)$ with $pb + r_o \alpha b - q > 0$
(positive expected return), \eqref{eq:kelly-foc} has a unique root
$w^\star \in (0, 1)$, and $\ell(w^\star) > 0$.
\end{proposition}

\begin{proof}
$\ell(w)$ is strictly concave on $(0, 1)$ since each summand is concave
in $w$. We have $\ell(0) = 0$ and $\ell'(0) = pb + r_o \alpha b - q$, which
is positive by hypothesis. Strict concavity guarantees a unique interior
maximum and $\ell(w^\star) > \ell(0) = 0$.
\end{proof}

We restrict the candidate menu to growth-positive spreads ($E > 0$);
spreads with $E \le 0$ are dropped before scoring. In practice
$E \le 0$ candidates rarely survive the credit-ratio gate of
Section~\ref{sec:methodology}.

\subsection{The entropic discount}

The Kullback--Leibler divergence between two probability distributions
$P$ and $Q$ on the same three-outcome space is
\begin{equation}
\KL(P \,\|\, Q) \;=\; \sum_{i \in \{p, r_o, q\}} P_i \log \frac{P_i}{Q_i}.
\end{equation}
$\KL$ is nonnegative and zero iff $P = Q$ almost everywhere. The
entropic discount factor $\exp(-k \cdot \KL) \in (0, 1]$ multiplies the
Kelly EV by a number that equals 1 at perfect agreement and shrinks
as $P$ and $Q$ diverge.

\paragraph{The choice of $P$ and $Q$.}
This is the single most consequential modelling decision in GROUND.
The original \citet{mercurio2020} paper used $Q = u$, the uniform
distribution, with $P = $ the candidate's empirical (or delta-induced)
distribution; high divergence then signals \emph{concentrated} bets,
which the discount penalizes. Later canonicals replaced $u$ with an
empirical bucket-conditioned reference $P_{\mathrm{hist}}$ estimated
from historical realizations of similarly-bucketed spreads. The
canonical of this paper introduces a third option developed in
Section~\ref{sec:rvvsiv-choice}: $P$ and $Q$ are both derived from
the BS $d_2$ on the candidate's strikes, one using the underlying's
realized vol and the other using the option's implied vol per leg.
The discount thus measures \emph{per-spread variance risk premium
expressed in probability space} rather than in vol space.

\subsection{The Hansen--Sargent identity}

\begin{proposition}[GROUND-Hansen-Sargent identity]
\label{prop:hs}
On the growth-positive cone $\{E(p, x) > 0\}$, the intrinsic GROUND
ranker is exactly the exponential of the Hansen-Sargent multiplier-preferences
functional:
\begin{equation}
\ln \GROUND(p, x, k) \;=\; g(p, x) - k \cdot \KL(P \,\|\, Q) \;=\; J_k(p, x),
\label{eq:hs-identity}
\end{equation}
where $g := \ln E$.
\end{proposition}

\begin{proof}
Take logarithms of \eqref{eq:ground}. The first factor gives $\ln E = g$;
the second factor gives $-k \cdot \KL$.
\end{proof}

\begin{remark}
\eqref{eq:hs-identity} is not an approximation; it is an exact identity
by construction. The multiplicative form \eqref{eq:ground} and the
additive form \eqref{eq:hs} are equivalent rankers since $\exp$ is
strictly monotone. The advantage of the multiplicative form is
interpretability: $\GROUND$ reads directly as a per-trade
risk-adjusted expected return. The additive form is the more familiar
one in the decision-theoretic literature.
\end{remark}


\section{The rv-vs-iv divergence choice}
\label{sec:rvvsiv-choice}

\subsection{Motivation: per-spread VRP in probability space}

For a credit spread, the seller is paid an implied-volatility premium
relative to the volatility actually realized over the spread's life.
This is the per-spread variance risk premium (VRP). At the
underlying-level, VRP is well documented: VIX exceeds 30-day realized vol
on average by 3--5 vol points \citep{bollerslev2009}. At the
spread-level the corresponding object is the divergence between the
probability distribution $Q_{\mathrm{IV}}$ induced by the option's
implied vol and the distribution $P_{\mathrm{RV}}$ induced by the
underlying's realized vol over a tenor-matched window. We make this
precise.

\paragraph{Construction.}
Given the spread's strikes $K_{\text{short}}, K_{\text{long}}$,
underlying spot $S$, leg implied volatilities $\sigma^{\mathrm{IV}}_{s},
\sigma^{\mathrm{IV}}_{\ell}$, days to expiry $T$, and a realized
volatility estimate $\sigma^{\mathrm{RV}}$ (computed over a window
tenor-matched to the option's life, Section~\ref{sec:methodology}), we
define
\begin{equation}
d_2(\sigma, K) := \frac{\ln(S/K) - \tfrac{1}{2} \sigma^2 T}{\sigma \sqrt{T}},
\end{equation}
and define the per-leg risk-neutral ITM probability under each volatility
hypothesis. For puts, $\Pr[S_T < K] = \Phi(-d_2)$; for calls,
$\Pr[S_T > K] = \Phi(d_2)$. Combining the two legs gives a three-outcome
probability vector $(p, r_o, q)$ on $\{$full credit kept, partial,
full loss$\}$ under each hypothesis:
\begin{equation}
P_{\mathrm{RV}} := (p^{\mathrm{RV}}, r_o^{\mathrm{RV}}, q^{\mathrm{RV}}), \qquad
Q_{\mathrm{IV}} := (p^{\mathrm{IV}}, r_o^{\mathrm{IV}}, q^{\mathrm{IV}}).
\end{equation}
The candidate's entropic discount is the forward KL divergence of these
two distributions:
\begin{equation}
\KL(P_{\mathrm{RV}} \,\|\, Q_{\mathrm{IV}}) \;=\;
\sum_{i \in \{p, r_o, q\}} P^{\mathrm{RV}}_i \log \frac{P^{\mathrm{RV}}_i}{Q^{\mathrm{IV}}_i}.
\end{equation}
This is dominated by the win-leg disagreement (the term weighted by
$p^{\mathrm{RV}}$ which is typically the largest mass under our menu
of OTM spreads). Concretely: if realized vol is much smaller than
implied vol, $\sigma^{\mathrm{IV}} \gg \sigma^{\mathrm{RV}}$, then
$\Phi(\cdot)$ at the short strike is much closer to 1 under the
$P_{\mathrm{RV}}$ measure than under $Q_{\mathrm{IV}}$; the KL term
becomes large and the discount factor shrinks.

\subsection{Tenor matching}

The realized-vol estimator $\sigma^{\mathrm{RV}}$ is a rolling-window
standard deviation of daily log returns, annualized by $\sqrt{252}$.
The window length \emph{must} be tenor-matched to the option's life.
Our weekly menu has $\mathrm{DTE} \in \{1, 2, 3, 4\}$; we use a
10-trading-day rolling window. A 30-day window was tested and
empirically failed (Section~\ref{sec:sensitivity-rv-window}): the
tenor mismatch causes the discount to be dominated by short-horizon
vol-regime memory that does not actually inform the next 1--4 day
outcome.

\subsection{Head-to-head against the empirical-bucket reference}
\label{sec:rvvsiv-vs-empirical}

We compared the rv\_vs\_iv choice head-to-head against the prior
empirical-bucket-conditioned reference (where $P$ is estimated by
binning historical realizations into $(\text{DTE}, \text{putcall},
\text{delta-bucket}, \text{IV-bucket}, \text{IV-rank-bucket})$
cells and querying the cell-realization frequency for the candidate's
bucket). Both rankers used the same Hansen-Sargent identity
\eqref{eq:hs-identity}, the same Kelly EV calculation, and the same
candidate menu. Both had directional gating disabled (no regime gate).
Both used the same Mon-Thu entry, Friday expiry, top-5 per
day-of-week selection with a per-DOW GROUND threshold. The amplification
parameter $k$ was tuned separately for each reference choice by a
two-dimensional sweep over $(k, \mathrm{threshold})$. Headline weekly
Sharpe on the 2020--2025 window:
\begin{center}
\begin{tabular}{lrrr}
\toprule
Divergence reference & Sharpe & Drawdown & Picks \\
\midrule
Empirical bucket ($P_{\mathrm{hist}} \,\|\, Q_{\mathrm{delta}}$) & 2.04 & $-8.8\%$ & 1{,}142 \\
rv\_vs\_iv ($P_{\mathrm{RV}} \,\|\, Q_{\mathrm{IV}}$) $\bigstar$  & 2.15 & $-7.4\%$ & 1{,}111 \\
\bottomrule
\end{tabular}
\end{center}
The rv\_vs\_iv reference wins on every metric. The reasons are
structural rather than parametric: (i) the empirical-bucket reference
introduces sampling noise tied to small cell counts (especially in
the joint $(\text{DTE}, \text{delta-bucket}, \text{IV-bucket}, \text{IV-rank-bucket})$
key), (ii) the rv\_vs\_iv reference captures a forward-looking risk
signal that bucket history can not, namely the disagreement between
the market's current pricing and recent realized vol.

\subsection{Why this works empirically}
\label{sec:rvvsiv-works}

The empirical pattern we observe across 22{,}702 pre-filter
candidates is striking and deserves explicit discussion.

\begin{enumerate}
\item \textbf{The divergence in isolation has no predictive power.}
Pearson correlation between $\KL$ and realized P\&L across all
candidates is $+0.004$. Within decile bins of $\KL$, mean realized
P\&L is negative in all 10 deciles (it has no direct $\KL \to $ P\&L
gradient).
\item \textbf{The divergence is a tail filter on top-$g$ candidates.}
Comparing the top-542 selection by $g$ alone against the top-542
selection by GROUND, GROUND retains 370 of the 542 (68\%) and swaps
out 172 (32\%). The 172 picks GROUND drops have mean P\&L $-\$1.26$
and win-rate 27.3\%; the 172 picks it adds have mean P\&L $+\$7.16$
and win-rate 51.7\%. The swap is uniform-in-WR: dropped picks have
median $\KL = 0.284$ (large), added picks have median $\KL = 0.005$
(near zero). The interpretation is that high-$g$ candidates with high
$\KL$ are precisely those where the market's implied vol is signaling
something that recent realized vol does not see (earnings, news, or
regime shift). The discount filters these out.
\item \textbf{The improvement is real, not random.} On the
full-sample 2020--2025 evaluation, GROUND top-542 produces mean
$+\$24.77$ per pick at 52.6\% win rate; $g$ alone produces mean
$+\$22.10$ at 44.8\%. The 7.8 percentage-point win-rate improvement
at unchanged pick count is what drives the Sharpe lift from $g$ alone
($\sim$1.5) to GROUND ($\sim$2.0).
\end{enumerate}

The full proof of this swap is given in Section~\ref{sec:proof}.


\section{Methodology}
\label{sec:methodology}

\subsection{Universe and candidate construction}

The universe is the SP100 set of largest-capitalization US equities
underlying liquid weekly options. Every Mon--Thu we screen for
single-strike vertical credit spreads expiring on Friday of the
current week: bull-puts (sell higher-strike put, buy adjacent
lower-strike put) and bear-calls (sell lower-strike call, buy adjacent
higher-strike call). Candidate filters:
\begin{itemize}
\item Days to expiry $\in \{1, 2, 3, 4\}$.
\item Short-leg delta $|\Delta| \in [0.35, 0.65]$ (we target near-ATM
exposure where credit-to-risk is highest while keeping the partial
zone narrow).
\item Both legs carry open interest $\geq 100$.
\item Spread width is one strike interval (we do not widen the strike
pair; wider spreads were tested and produced lower per-spread Sharpe).
\item Credit-to-risk ratio $b \geq 0.30$ and $\leq 5$ (very-rich-credit
spreads correspond to deep-ITM short legs that are not the intended
strategy).
\end{itemize}

\subsection{Entry credit basis}
\label{sec:entry-credit}

The entry credit used in backtest is
\[
c_{\mathrm{entry}} = 0.80 \times c_{\mathrm{LAST}},
\]
where $c_{\mathrm{LAST}} = \max(\mathrm{bid}_s, \min(\mathrm{LAST}_s, \mathrm{ask}_s))
- \max(\mathrm{bid}_\ell, \min(\mathrm{LAST}_\ell, \mathrm{ask}_\ell))$ is
the clamped-LAST credit per share. The 0.80 multiplier is a
conservative 20\% haircut to account for the realistic difference
between modelled close-LAST and the price at which a limit order is
actually filled; sensitivity to this choice is given in
Section~\ref{sec:sensitivity-haircut}.

\subsection{No regime gate}
\label{sec:no-regime}

Prior canonicals applied a directional gate that allowed bull-puts
only when SPY closed above its 100-day SMA (``bull regime'') and
bear-calls only below it. The 2026-06 canonical drops this gate. The
empirical justification is straightforward: in the 2020--2025 sample,
of the top-quintile-GROUND picks, the bull-puts placed in bear regimes
have a 58\% win rate and the bear-calls placed in bull regimes
contribute $+\$11{,}000$ to total P\&L. The regime gate was simply
filtering out edge that the GROUND ranker had already identified.

\subsection{Earnings and ex-dividend gates}
\label{sec:gates}

In live trading we additionally filter the candidate set with two
data-driven gates: (i) earnings gate, drop any candidate whose
underlying has an earnings announcement scheduled within the
$[\mathrm{entry}, \mathrm{expiry}]$ window (both bull-puts and
bear-calls); (ii) ex-dividend gate, drop any \emph{bear-call} whose
underlying has an ex-dividend date within
$[\mathrm{entry}, \mathrm{expiry} + 1]$ (only bear-calls; bull-puts
are not at risk of early assignment for dividend capture). These
gates address tail risks that are not in the in-sample backtest
because announcement dates are not historically encoded in the option
data. In live operation the gates are populated daily from NASDAQ's
public calendar API.

\subsection{Partial-zone realization}
\label{sec:partial-zone}

In live operation, a spread closing in the partial zone on Friday
exposes the position to weekend assignment risk on the short leg.
Avoiding assignment requires closing the position prior to 4pm
settlement, which on illiquid weekly options carries a fill-quality
penalty above intrinsic value. The 2020--2025 backtest reports both
the raw piecewise-linear payoff (which assumes settlement at close)
and a \emph{realistic} variant that haircuts partial-WIN outcomes to
50\% of their intrinsic value (modelling the typical pin-risk
close-cost) while keeping partial-LOSS outcomes at 100\% (which
cannot be improved through early closing). Both headlines are
reported.


\section{Headline results: 2020--2025}
\label{sec:results}

\subsection{Aggregate performance}

Table~\ref{tab:headline} reports the headline statistics for the
canonical 2020--2025 backtest under fixed-qty=2 sizing, qty=1 sizing,
and a fractional-Kelly sizing rule of $\tfrac{1}{16}$ of the per-trade
Kelly fraction with a per-spread contract cap of 5. We report both the
raw payoff variant (partial wins at full intrinsic) and the
pin-risk-realistic variant (partial wins at 50\% of intrinsic).

\begin{table}[h]
\centering
\caption{Headline backtest, 2020-01-07 through 2025-12-31 (\textbf{1{,}111 trades}, 297 weeks).}
\label{tab:headline}
\begin{tabular}{lrrrr}
\toprule
Sizing & Final \$ & Total Return & Sharpe (weekly) & MaxDD \\
\midrule
\multicolumn{5}{l}{\emph{Raw payoff (partial wins at full intrinsic):}}\\
qty = 1                         & 36{,}906 & 269.1\%  & 2.19 & $-5.6\%$ \\
qty = 2 (canonical)             & 63{,}812 & 538.1\%  & 2.15 & $-7.4\%$ \\
$\tfrac{1}{16}$-Kelly (cap = 5) & 78{,}801 & 688.0\%  & 2.24 & $-5.3\%$ \\
\midrule
\multicolumn{5}{l}{\emph{Pin-risk-realistic (partial wins at 50\% intrinsic):}}\\
qty = 1                         & 34{,}566 & 245.7\%  & 2.08 & $-6.1\%$ \\
qty = 2 (canonical)             & 59{,}131 & 491.3\%  & 2.05 & $-8.0\%$ \\
$\tfrac{1}{16}$-Kelly (cap = 5) & 72{,}382 & 623.8\%  & 2.17 & $-5.7\%$ \\
\midrule
SPY buy-and-hold (same window)  & 21{,}018 & 110.2\%  & 0.71 & $-34.1\%$ \\
\bottomrule
\end{tabular}
\end{table}

\paragraph{Reading the table.}
Under the canonical fixed-qty=2 sizing rule with the pin-risk haircut,
\$10{,}000 starting capital compounds to \$59{,}131, an annualized
return of $33.4\%$ over the 5.97-year window. The weekly Sharpe of
$2.05$ is well above what SPY buy-and-hold delivered over the same
window ($0.71$), and the maximum drawdown of $8.0\%$ is approximately
one-fourth of SPY's.

\paragraph{Direction breakdown.}
Of the 1{,}111 picks in the canonical 2020--2025 backtest, 593 (53.4\%)
are bear-calls and 518 (46.6\%) are bull-puts. Bear-calls contributed
the majority of total P\&L (\$32{,}410 vs.\ \$16{,}721 for bull-puts);
the per-pick mean P\&L is $+\$54.70$ for bear-calls vs.\ $+\$32.30$
for bull-puts. Win rates are within 1pp of each other (54.1\% vs.\
53.3\%) so the directional asymmetry is not a hit-rate effect; the
bear-call advantage is a credit-richness effect driven by persistent
call-side skew premium in the underlying market.


\section{Out-of-time validation: 2026}
\label{sec:oot}

The 2026 holdout is a clean out-of-time test in the sense that no
parameter selection or methodology choice in this paper references
2026 data. The $(k, \mathrm{threshold})$ pair was selected on the
2020--2025 sample; all qualitative methodology choices (rv\_vs\_iv
divergence, tenor-matched 10-day RV, no regime gate, earnings and
ex-div gates, partial-WIN haircut) were settled before 2026 was scored.

\paragraph{Procedure.}
For each Mon--Thu in 2026 H1 we scored the SP100 candidate menu under
the canonical settings and computed realized P\&L from the vendor's
end-of-Friday underlying closing print. Sizing is fixed qty = 2 to
match the in-sample headline.

\begin{table}[h]
\centering
\caption{Out-of-time validation, 2026-01 through 2026-06.}
\label{tab:oot}
\begin{tabular}{lrrr}
\toprule
Statistic & Value & In-sample 2020--2025 & SPY (same window) \\
\midrule
Picks                     & [OOT-PICKS]   & 1{,}111 & --- \\
Weeks                     & [OOT-WEEKS]   & 297     & --- \\
Final \$ (from \$10{,}000) & [OOT-FINAL]  & 59{,}131 & [OOT-SPY-FINAL] \\
Total return              & [OOT-RET]\%   & 491.3\%  & [OOT-SPY-RET]\% \\
Sharpe (weekly)           & [OOT-SHARPE]  & 2.05     & [OOT-SPY-SHARPE] \\
Max drawdown              & [OOT-DD]\%    & $-8.0\%$ & [OOT-SPY-DD]\% \\
Win rate                  & [OOT-WR]\%    & 53.7\%   & --- \\
\bottomrule
\end{tabular}
\end{table}

[OOT-DISCUSSION-PARAGRAPH]


\section{GROUND vs.\ \emph{g} alone: the swap proof}
\label{sec:proof}

The empirical content of GROUND beyond Kelly EV is the entropic
discount $\exp(-k \cdot \KL)$. This section establishes by direct
comparison that the discount adds material selection value on the
realized sample, despite the fact that $\KL$ has no detectable direct
correlation with P\&L outcome.

\subsection{Same-N comparison}

Across the 22{,}702 pre-filter candidates of the 2020--2025 universe,
we form three top-$N$ selection rules and compare their realized
outcomes:
\begin{enumerate}
\item Top-542 by $g$ alone (Kelly EV, no discount).
\item Top-542 by $-\KL$ alone (lowest divergence, no Kelly EV).
\item Top-542 by GROUND ($\bigstar$).
\end{enumerate}
542 is the count above the canonical GROUND $\geq 0.075$ threshold.

\begin{table}[h]
\centering
\caption{Same-$N$ head-to-head, top-542 selection across rules.}
\label{tab:swap-head}
\begin{tabular}{lrrrr}
\toprule
Selection rule & $n$ & Mean P\&L & Sum P\&L & Win rate \\
\midrule
$g$ alone (top 542)               & 542 & $+\$22.10$ & $+\$11{,}975$ & 44.8\% \\
$-\KL$ alone (bottom 542 by $\KL$) & 542 & $-\$20.00$ & $-\$10{,}841$ & 48.0\% \\
GROUND top 542 $\bigstar$          & 542 & $+\$24.77$ & $+\$13{,}424$ & 52.6\% \\
\bottomrule
\end{tabular}
\end{table}

GROUND adds $\$2.67$ per pick over $g$ alone at unchanged pick count,
and lifts the win rate by 7.8 percentage points. Ranking by $-\KL$
alone is a money loser: spreads with $P \approx Q$ (low divergence)
have no special edge by themselves.

\subsection{The 32\% swap}

GROUND's top-542 retains 370 of the 542 top-$g$ candidates (68\%) and
swaps out 172 (32\%). The swap-in and swap-out subsets have very
different realized outcomes:

\begin{table}[h]
\centering
\caption{The 172-pick swap: which picks GROUND drops vs.\ adds.}
\label{tab:swap}
\begin{tabular}{lrrrr}
\toprule
Subset & $n$ & Mean P\&L & Win rate & Median $\KL$ \\
\midrule
In $g$-top, NOT in GROUND-top (\emph{dropped by} $\KL$) & 172 & $-\$1.26$ & 27.3\% & 0.2838 \\
In GROUND-top, NOT in $g$-top (\emph{added by} $\KL$)   & 172 & $+\$7.16$ & 51.7\% & 0.0052 \\
\bottomrule
\end{tabular}
\end{table}

The dropped subset has a win rate of \emph{27.3\%}: roughly one in
four. Their median $\KL$ is 0.28, i.e., the IV-implied and RV-implied
distributions disagree substantially. The added subset has a win rate
of 51.7\% and median $\KL$ near zero (IV and RV agree). The swap
captures roughly $\$1{,}450$ of edge across the 5.5-year sample.

\subsection{Interpretation}

The empirical pattern is that $\KL$ acts as a \emph{tail filter} on
Kelly EV, not a per-candidate predictor. Its Pearson correlation with
realized P\&L across all 22{,}702 candidates is $+0.004$; in decile
analysis it produces no monotone gradient on its own. But applied to
the top-$g$ tail, the discount swaps 27\%-win-rate candidates for
52\%-win-rate candidates, lifting the per-pick mean by 12\% and the
win rate by 7.8 percentage points.

The structural interpretation is that high-$\KL$ candidates within
the top-$g$ tail are precisely those for which the market's implied
vol is signaling a risk realized vol cannot see. The discount
implements a form of ``trust the market when it disagrees with
history'' that complements the ``trust history when it has edge''
content of $g$.


\section{Sensitivity analyses}
\label{sec:sensitivity}

\subsection{Entry credit haircut}
\label{sec:sensitivity-haircut}

The entry credit basis is the single most consequential modelling
choice for live deployability of the strategy. Below we sweep the
$c_{\mathrm{entry}} / c_{\mathrm{LAST}}$ multiplier from 0.95 to 0.50
on the canonical 1{,}111-pick cache and report the realized payoff:

\begin{table}[h]
\centering
\caption{Entry-credit-basis sensitivity (qty = 2).}
\label{tab:credit-sweep}
\begin{tabular}{lrrrr}
\toprule
Basis & Mean credit & Mean P\&L & Final \$ & Sharpe-W \\
\midrule
$0.95 \times$ LAST  & \$1.43 & $+\$46.74$ & 103{,}851 & 2.61 \\
$0.90 \times$ LAST  & \$1.35 & $+\$39.23$ & 87{,}171  & 2.42 \\
$0.85 \times$ LAST  & \$1.28 & $+\$31.72$ & 70{,}492  & 2.19 \\
$0.80 \times$ LAST ($\bigstar$) & \$1.20 & $+\$24.22$ & 53{,}812 & 1.93 \\
$0.75 \times$ LAST  & \$1.13 & $+\$16.71$ & 37{,}133  & 1.62 \\
$0.70 \times$ LAST  & \$1.05 & $+\$9.21$  & 20{,}454  & 1.20 \\
$0.65 \times$ LAST  & \$0.98 & $+\$1.70$  & 3{,}774   & 0.39 \\
$0.60 \times$ LAST  & \$0.90 & $-\$5.81$  & $-12{,}905$ & $-1.07$ \\
\bottomrule
\end{tabular}
\end{table}

The strategy breaks even at approximately $0.65 \times$ LAST, i.e., a
35\% haircut on the last-traded mark. The canonical $0.80$ multiplier
leaves a 15-point buffer above breakeven, sized to accommodate the
realistic limit-fill-quality gap in live trading.

\subsection{Partial-WIN haircut}
\label{sec:sensitivity-partialwin}

Section~\ref{sec:partial-zone} introduced the canonical 50\%
intrinsic haircut on partial WINS. Pure raw-payoff and pure
worst-case (partial-WIN $\to$ partial-LOSS = max-loss) bracket the
realistic case:

\begin{table}[h]
\centering
\caption{Partial-WIN realization sensitivity (qty = 2, $0.80 \times$ LAST entry).}
\label{tab:partial-sweep}
\begin{tabular}{lrrrr}
\toprule
Partial-WIN treatment & Final \$ & Sharpe-W & MaxDD & Win rate \\
\midrule
At full intrinsic (raw)       & 63{,}812 & 2.15 & $-7.4\%$ & 53.7\% \\
At 50\% intrinsic ($\bigstar$) & 59{,}131 & 2.05 & $-8.0\%$ & 53.7\% \\
At 0\% (treat as breakeven)   & 56{,}450 & 1.99 & $-8.6\%$ & 53.7\% \\
At $-100\%$ (max loss, worst case) & 36{,}940 & 1.39 & $-15.2\%$ & 45.7\% \\
\bottomrule
\end{tabular}
\end{table}

Even under the brutal worst case (every partial closes at max loss),
the strategy still produces Sharpe 1.39 and a manageable 15.2\%
drawdown, more than twice SPY's risk-adjusted return on the same
window. The 50\%-intrinsic canonical sits comfortably between the raw
and worst-case extremes and is consistent with anecdotal live
close-cost data on illiquid OTM weekly spreads.

\subsection{Probability-partition collapses}

\noindent We tested two collapses of the three-outcome partition into a
two-outcome one and reported the impact on canonical headline metrics:

\begin{table}[h]
\centering
\caption{Probability-partition sensitivity.}
\label{tab:partition-sweep}
\begin{tabular}{lrrrr}
\toprule
Partition & Picks & Sharpe-W & MaxDD & Final \$ \\
\midrule
3-state, $\alpha$-weighted ($\bigstar$) & 1{,}111 & 2.05 & $-8.0\%$ & 59{,}131 \\
2-state, fold $r_o$ into $q$ (partial $\to$ loss) & ~700 & 1.89 & $-13.6\%$ & 52{,}470 \\
2-state, fold $r_o$ into $p$ (partial $\to$ win) & ~3{,}700 & 0.65 & $-73.6\%$ & 26{,}206 \\
\bottomrule
\end{tabular}
\end{table}

Both partition-collapse variants underperform the canonical 3-state.
The ``partial $\to$ loss'' variant selects too few, too-deep-OTM
spreads that lose harder when they lose; the ``partial $\to$ win''
variant is grossly over-permissive and catastrophic. The
$\alpha$-weighted partial-zone multiplier (Section~\ref{sec:framework})
is doing real work; it is not a notation choice.

\subsection{RV window length}
\label{sec:sensitivity-rv-window}

We tested $\sigma^{\mathrm{RV}}$ computed over a 30-calendar-day
trailing window (21 trading days), as a tenor mismatch against the
$\mathrm{DTE} \in \{1, 2, 3, 4\}$ option menu. The strategy still
profited but Sharpe dropped to $1.49$ from the tenor-matched
10-day-window canonical's $2.05$. The 10-day window was selected as
the shortest window with $\geq 5$ daily-return observations supporting
a stable standard-deviation estimate.

\subsection{Regime gate (off vs.\ on)}

The 2026-06 canonical drops the binary directional regime gate that
prior canonicals applied. Re-introducing it (bull-puts only above
SPY's 100-day SMA, bear-calls only below) reduces picks from 1{,}111
to 541 and reduces weekly Sharpe from $2.05$ to $1.48$. The gate
filters out, e.g., bear-call placements in bull regimes that
contribute the majority of bear-call P\&L. The unconditional ranker
identifies edge that the gate is hiding.


\section{Conclusion}
\label{sec:conclusion}

The intrinsic GROUND ranker, applied to SP100 weekly credit spreads
with the rv\_vs\_iv divergence reference described in this paper,
produces a weekly Sharpe of $2.05$ on a 5.97-year in-sample window
and weekly Sharpe of \textbf{[OOT-SHARPE]} on a strict 2026 holdout.
Two structural contributions distinguish this canonical from prior
GROUND deployments. First, the divergence reference is built from
realized and implied volatility through the Black-Scholes $d_2$,
producing a per-spread variance-risk-premium signal directly in
probability space. Second, the binary directional regime gate is
removed: the unconditional GROUND ranker discovers per-spread edge
in both bull-put and bear-call directions across both market regimes,
and the regime gate was filtering out true edge.

The discount factor's role in selection is subtle: across the
22{,}702 pre-filter candidates of the in-sample window, the
divergence's direct correlation with realized P\&L is $+0.004$, but
applied as a multiplicative discount to Kelly EV at the top-$g$ tail
it swaps 27\%-win-rate candidates for 52\%-win-rate ones, lifting
the per-pick mean by 12\% and the win rate by 7.8 percentage points.
The mechanism is interpretable: high-$\KL$ candidates within the
top-$g$ tail are spreads where the market's implied vol is
signaling something realized vol does not see; the discount filters
these out.

We close by emphasizing what the GROUND ranker is and is not. It is
a top-tail skimmer; nine of ten deciles of the candidate distribution
are not-EV-positive on average. The selection threshold is not
optional, it is the strategy. The 2026 holdout result confirms that
the parameters fixed on 2020--2025 transfer cleanly to a held-out
window with no evidence of in-sample over-fitting.


\section*{Acknowledgements}

The author thanks the open-source maintainers of \texttt{pandas},
\texttt{numpy}, and \texttt{ib\_insync} for the infrastructure used in
the live implementation, and the IBKR options market data feed
through which the live rankings are produced.


\bibliographystyle{plainnat}
\begin{thebibliography}{99}

\bibitem[Bollerslev et~al.(2009)]{bollerslev2009}
T.~Bollerslev, G.~Tauchen, and H.~Zhou.
\newblock Expected stock returns and variance risk premia.
\newblock \emph{Review of Financial Studies}, 22(11):4463--4492, 2009.

\bibitem[Hansen and Sargent(2008)]{hansen2008robustness}
L.\,P.~Hansen and T.\,J.~Sargent.
\newblock \emph{Robustness}.
\newblock Princeton University Press, 2008.

\bibitem[Israelov(2017)]{israelov2017}
R.~Israelov.
\newblock Pathetic protection: The elusive benefits of protective puts.
\newblock \emph{Journal of Alternative Investments}, 19(3):38--56, 2017.

\bibitem[Lo(2002)]{lo2002}
A.\,W.~Lo.
\newblock The statistics of Sharpe ratios.
\newblock \emph{Financial Analysts Journal}, 58(4):36--52, 2002.

\bibitem[Maccheroni et~al.(2006)]{maccheroni2006ambiguity}
F.~Maccheroni, M.~Marinacci, and A.~Rustichini.
\newblock Ambiguity aversion, robustness, and the variational
representation of preferences.
\newblock \emph{Econometrica}, 74(6):1447--1498, 2006.

\bibitem[Mercurio(2020)]{mercurio2020}
P.\,J.~Mercurio.
\newblock The GROUND ratio: A portfolio ranking criterion combining
growth and entropic risk.
\newblock Working paper, 2020.

\bibitem[Mercurio(2020b)]{mercurio2020portfolio}
P.\,J.~Mercurio.
\newblock Entropy as a risk measure in portfolio optimization.
\newblock Working paper, 2020.

\bibitem[Mertens(2002)]{mertens2002}
E.~Mertens.
\newblock Comments on variance of the IID estimator in Lo (2002).
\newblock Unpublished working paper.

\bibitem[Opdyke(2007)]{opdyke2007}
J.\,D.~Opdyke.
\newblock Comparing Sharpe ratios: So where are the $p$-values?
\newblock \emph{Journal of Asset Management}, 8(5):308--336, 2007.

\end{thebibliography}

\end{document}
